\documentclass[11pt,a4paper]{article}
\usepackage[T1]{fontenc}
\usepackage[utf8]{inputenc}
\usepackage[main=italian,provide=*]{babel}
\usepackage{amsmath,amssymb}
\usepackage{geometry}
\geometry{margin=2.5cm}
\usepackage{booktabs}
\usepackage{seqsplit}
\usepackage{hyperref}
\hypersetup{colorlinks=true,linkcolor=blue,citecolor=blue,urlcolor=blue}

\title{Manuale d'uso --- correlazioni dinamiche trimero a catena aperta\\
\large versione attuale (circuito VQE reale \texttt{W-2qC.K2})}
\author{Note di lavoro --- Tesi triennale, Samuele Schiavi\\ Relatore: Prof.\ Alessandro Chiesa}
\date{}

\begin{document}
\maketitle

\section*{Scopo}

Questo manuale spiega come eseguire la simulazione delle correlazioni dinamiche
$C_{ij}^{\alpha\beta}(t)$ per il trimero a catena aperta nella versione \emph{attuale}: lo
stato fondamentale non è più preparato con le ampiezze esatte, ma con il circuito VQE reale
(ansatz \texttt{W-2qC.K2}, $10$ parametri, $2$ cicli), ottimizzato numericamente. La versione a
stato esatto resta disponibile (comportamento di default se non si passano parametri VQE) e
serve da riferimento: vedi il manuale gemello
\texttt{\seqsplit{manuale\_uso\_correlazioni\_trimero\_catena\_statevector.tex}}.

\section*{Prerequisiti}

Tutti i file del manuale gemello, più:
\begin{itemize}
\item \texttt{vqe\_w2qC\_k2\_trimero\_catena.py} --- costruzione dell'ansatz \texttt{W-2qC.K2}
e ottimizzazione VQE, su \textbf{entrambi} i punti di lavoro.
\item \texttt{validate\_vqe\_circuito\_correlazioni\_catena.py} --- confronto sistematico contro
la versione a stato esatto, entrambi i punti.
\end{itemize}

\begin{quote}
\textbf{Importante, a differenza dell'anello}: esistono \textbf{due} file di parametri
ottimali, uno per punto di lavoro --- \texttt{w2qC\_k2\_params\_vqedm.npz} (VQE-DM,
$J{=}1,b{=}3.0,D{=}0.15$) e \texttt{w2qC\_k2\_params\_S0.npz} (S0,
$J{=}1,b{=}3.0073414,D{=}0.3$). Come nell'anello, i parametri dipendono dal punto: non riusare
quelli di un punto per l'altro.
\end{quote}

\section*{Passo 1 --- ottimizzazione VQE (una tantum per punto di lavoro)}

\begin{verbatim}
python vqe_w2qC_k2_trimero_catena.py
\end{verbatim}
Esegue $60$ partenze casuali (COBYLA) più polish (L-BFGS-B) \textbf{su entrambi i punti in
sequenza} (funzione \texttt{main()}), e salva due file. Valori ottenuti:
\begin{center}
\begin{tabular}{lcc}
\toprule
& VQE-DM & S0 \\
\midrule
$E_{\rm vqe}$ & $-7.36924699164432$ & $-7.75658090425334$ \\
$E_{\rm esatto}$ & $-7.36924699164433$ & $-7.75658090425337$ \\
$\mathcal F$ & $1.000000000000$ & $0.999999999999992$ \\
\bottomrule
\end{tabular}
\end{center}
Nella sessione di sviluppo sono bastati $15$ restart per raggiungere questi valori (verificato:
$5$ restart già sufficienti al punto VQE-DM); lo script di produzione ne usa $60$ per
sicurezza, seguendo la stessa cautela già documentata in
\texttt{\seqsplit{trimero\_catena\_vqe\_dm.tex}} per il punto $b=b_c$ (vicino a un level
crossing a $D=0$, dove pochi restart possono dare un falso ceiling). Se il tuo output mostra
$1-\mathcal F$ molto più grande (es.\ $10^{-3}$ o peggio), aumenta \texttt{n\_restarts} nella
chiamata a \texttt{optimize\_point}.

\section*{Passo 2 --- un correlatore singolo, con preparazione VQE}

\begin{verbatim}
import numpy as np
from circuito_correlazioni_trimero_catena import correlator_from_circuit

data = np.load("w2qC_k2_params_vqedm.npz")
params = data["params"]

J, b, D = 1.0, 3.0, 0.15
C = correlator_from_circuit(1, 'y', 1, 'y', t=1.3, N=200, J=J, b=b, D=D,
                             ansatz_params=params)
print(C)
\end{verbatim}
Passando \texttt{ansatz\_params}, la preparazione usa \texttt{w2qC\_k2\_circuit()} (X + due
cicli di [2 blocchi \texttt{W}, 3 $R_y$]) invece di \texttt{prepare\_state}. Per il punto S0,
caricare \texttt{w2qC\_k2\_params\_S0.npz} e usare $b{=}3.0073414, D{=}0.3$.

\section*{Passo 3 --- validazione contro la versione a stato esatto}

\begin{verbatim}
python validate_vqe_circuito_correlazioni_catena.py
\end{verbatim}
Richiede \textbf{entrambi} i file \texttt{.npz} del Passo 1 (già presenti). Esegue
automaticamente entrambi i punti in sequenza, ricalcola le 81 combinazioni con preparazione VQE
e con preparazione esatta, stampa l'errore medio e massimo:
\begin{center}
\begin{tabular}{lcc}
\toprule
& VQE-DM & S0 \\
\midrule
errore medio & $1.9\times10^{-8}$ & $3.8\times10^{-8}$ \\
errore massimo & $5.9\times10^{-8}$ & $9.8\times10^{-8}$ \\
$\sqrt{1-\mathcal F}$ (riferimento) & $4.5\times10^{-8}$ & $9.1\times10^{-8}$ \\
\bottomrule
\end{tabular}
\end{center}
Entrambi coerenti con $\sqrt{1-\mathcal F}$, come atteso per un residuo di fedeltà a livello di
stato. Errori molto più grandi indicano parametri VQE non ottimizzati bene (torna al Passo 1).

\section*{Passo 4 --- scan completo (statevector + shot), preparazione VQE}

A differenza dell'anello, non esiste uno script \texttt{scan81\_vqe\_*} separato: lo scan
completo con preparazione VQE è incorporato nei notebook (vedi sotto), non in uno script
standalone. Per un confronto rapido da riga di comando, usare \texttt{correlator\_from\_circuit}
con \texttt{ansatz\_params} in un ciclo sulle $81$ combinazioni (Passo 2, ripetuto).

\section*{Come esplorare tutto insieme: il notebook}

\texttt{\seqsplit{circuito\_correlazioni\_trimero\_catena\_vqe.ipynb}} carica i parametri già
ottimizzati (Passo 1, non li ricalcola) e mostra, con la preparazione VQE al posto di
\texttt{prepare\_state}: misura diretta delle 81 combinazioni, heatmap $|C|$ d'insieme, vista
nel tempo su griglia $9\times9$ (Re e Im), selettore per ispezione singola, verifica della
relazione $P_{13}$ ($C_{11}=C_{33}$), spettro dietro un caso ricco/piatto, confronto con lo scan
a shot finiti precomputato (\texttt{scan81\_catena\_vqedm\_results.npz}, Passo 4 del manuale
gemello). Lavora solo al punto VQE-DM (hard-coded); per S0 vanno adattati i parametri
$(J,b,D)$ e ricaricato \texttt{w2qC\_k2\_params\_S0.npz}.

\section*{Cosa NON cambia rispetto alla versione a stato esatto}

\texttt{build\_correlator\_circuit} e \texttt{correlator\_from\_circuit} sono retrocompatibili:
se non si passa \texttt{ansatz\_params} (o si passa \texttt{None}), il comportamento è
identico, bit per bit, alla versione statevector --- tutto il codice esistente
(\texttt{scan81\_trimero\_catena.py}, \texttt{\seqsplit{validate\_shot\_noise\_trimero\_catena.py}},
i primi tre notebook della fase) continua a funzionare senza modifiche.

\section*{Problemi comuni}

\begin{itemize}
\item \texttt{FileNotFoundError} su uno dei due file \texttt{w2qC\_k2\_params\_*.npz} --- va
generato eseguendo \texttt{vqe\_w2qC\_k2\_trimero\_catena.py} (Passo 1); richiede alcuni minuti
per via dei $60$ restart $\times$ $2$ punti.
\item Fidelity bassa dopo l'ottimizzazione --- vedi nota al Passo 1 (minimo locale, più probabile
vicino a $b_c$); aumentare i restart di solito risolve.
\item Uso dei parametri sbagliati per il punto --- verificare sempre quale \texttt{.npz} si sta
caricando: i due punti hanno parametri ottimali diversi e non intercambiabili.
\item Import circolare se si prova a importare \texttt{ground\_state} da
\texttt{\seqsplit{circuito\_correlazioni\_trimero\_catena.py}} dentro
\texttt{\seqsplit{vqe\_w2qC\_k2\_trimero\_catena.py}} --- per questo il secondo modulo definisce
una propria copia locale \texttt{\_ground\_state}, stessa scelta già fatta per l'anello.
\end{itemize}

\section*{Riferimenti}

Derivazione dell'ansatz: \texttt{\seqsplit{trimero\_catena\_vqe\_dm.tex}},
\texttt{ansatz\_catena.py}. Documentazione teorica: \texttt{\seqsplit{simmetrie\_correlatori\_trimero\_catena.tex}},
\texttt{\seqsplit{circuito\_correlazioni\_trimero\_catena\_spiegato.tex}}. Catalogo completo
della fase: \texttt{\seqsplit{trimero\_catena\_correlazioni.tex}}.

\end{document}
